\chapter*{ประมวลการสอนและแผนการจัดการเรียนรู้}
\addcontentsline{toc}{chapter}{ประมวลการสอนและแผนการจัดการเรียนรู้}

\section*{คำอธิบายรายวิชา}
การศึกษาระบบหน่วยและการวัดตามมาตรฐานสากล การวิเคราะห์เวกเตอร์ จลนศาสตร์และการเคลื่อนที่ในหนึ่งมิติและสองมิติ กฎการเคลื่อนที่ของนิวตันและแรงโน้มถ่วง โมเมนตัมและการชน สมดุลกลและโมเมนต์ งานและพลังงาน กำลังและเครื่องกลอย่างง่าย กลศาสตร์ของไหลและปรากฏการณ์ความตึงผิว ปรากฏการณ์ความร้อนและการถ่ายโอนความร้อน กฎของอุณหพลศาสตร์ ไฟฟ้าสถิตและศักย์ไฟฟ้าที่เยื่อหุ้มเซลล์ ไฟฟ้ากระแสตรงและวงจรไฟฟ้า ทัศนศาสตร์ทางเรขาคณิตและการทำงานของกล้องจุลทรรศน์ กัมมันตภาพรังสีและฟิสิกส์นิวเคลียร์ ตลอดจนการประยุกต์ใช้ในวิทยาศาสตร์ชีวภาพ การเกษตร และสิ่งแวดล้อม ควบคู่กับการฝึกทักษะการวัดละเอียดในห้องปฏิบัติการ

\section*{วัตถุประสงค์การเรียนรู้}
เมื่อสิ้นสุดการศึกษารายวิชานี้ นักศึกษาสามารถ
\begin{enumerate}
    \item อธิบายมโนทัศน์และกฎทางฟิสิกส์พื้นฐานที่เกี่ยวข้องกับการเคลื่อนที่ แรง พลังงาน ความร้อน ไฟฟ้า และคลื่นแม่เหล็กไฟฟ้าได้อย่างถูกต้อง
    \item วิเคราะห์และแก้ปัญหาโจทย์เชิงฟิสิกส์โดยใช้แบบจำลองทางคณิตศาสตร์และเวกเตอร์ได้อย่างเป็นระบบ
    \item เชื่อมโยงและประยุกต์ใช้หลักการทางฟิสิกส์ในการอธิบายปรากฏการณ์ของสิ่งมีชีวิตและเทคโนโลยีทางชีวภาพและสิ่งแวดล้อม
    \item ใช้อุปกรณ์วัดละเอียดทางฟิสิกส์ เช่น เวอร์เนียร์แคลิปเปอร์ ไมโครมิเตอร์ และมิเตอร์ไฟฟ้าได้อย่างถูกต้อง แม่นยำ และปลอดภัย
    \item วิเคราะห์ความคลาดเคลื่อนของการทดลองและเขียนรายงานผลการทดลองทางวิทยาศาสตร์ตามมาตรฐานสากล
\end{enumerate}

\section*{โครงสร้างแผนการสอน 13 สัปดาห์}

\begin{table}[H]
\centering
\small
\begin{tabularx}{\textwidth}{c p{5.5cm} X}
\toprule
\textbf{สัปดาห์} & \textbf{หัวข้อบทเรียน} & \textbf{กิจกรรมและปฏิบัติการ} \\
\midrule
1 & บทที่ 1 ระบบหน่วยและการวัด & ปฏิบัติการเครื่องมือวัดละเอียด (เวอร์เนียร์/ไมโครมิเตอร์) \\
2 & บทที่ 2 เวกเตอร์ในทางฟิสิกส์ & การรวมแรงแบบเวกเตอร์บนโต๊ะทดลองแรง \\
3 & บทที่ 3 จลนศาสตร์และการเคลื่อนที่ & รางปล่อยเทปตรวจจับเวลาการเคลื่อนที่แนวดิ่ง \\
4 & บทที่ 4 กฎการเคลื่อนที่ของนิวตัน & แรงเสียดทานและเครื่องชั่งสปริง \\
5 & บทที่ 5 โมเมนตัมและการดล & การชนแบบยืดหยุ่นบนรางลมอัด \\
6 & บทที่ 6 สมดุลกลและโมเมนต์ & สภาพสมดุลของคานและจุดศูนย์ถ่วง \\
7 & บทที่ 7 งานและพลังงานกล & กฎการอนุรักษ์พลังงานในเพนดูลัม \\
8 & บทที่ 8 กำลังและเครื่องกลอย่างง่าย & ประสิทธิภาพเครื่องกล รอก และพื้นเอียง \\
9 & บทที่ 9 ปรากฏการณ์ความร้อน & คาโลริมิเตอร์และความร้อนจำเพาะของสสาร \\
10 & บทที่ 10 อุณหพลศาสตร์และก๊าซ & กฎของก๊าซในอุดมคติและเครื่องยนต์ความร้อน \\
11 & บทที่ 11 ไฟฟ้าสถิตและสนามไฟฟ้า & การกระจายประจุและศักย์ไฟฟ้า \\
12 & บทที่ 12 ไฟฟ้ากระแสตรงและวงจร & กฎของโอห์มและการต่อความต้านทาน \\
13 & บทที่ 13 ทัศนศาสตร์และกัมมันตภาพรังสี & เลนส์นูน เลนส์เว้า และการทำงานของกล้องจุลทรรศน์ \\
\bottomrule
\end{tabularx}
\caption{แผนการจัดการเรียนรู้รายสัปดาห์วิชาฟิสิกส์พื้นฐาน}
\end{table}
